\section{Conceptos fundamentales}


    \textbf{Gravimetría:} Método de análisis cuantitativo basado en la medición de la masa de un compuesto de composición química definida, relacionado con el analito.
    
    \textbf{Factor gravimétrico (FG):} Relación entre la masa molar del analito y la masa molar del compuesto pesado, multiplicada por la relación estequiométrica.
    
    \textbf{Precipitado:} Sólido formado a partir de una reacción química que contiene al analito.
    
    \textbf{Muestra seca:} Muestra a la que se le ha eliminado la humedad para obtener un peso constante.


\section{Tipos de análisis gravimétrico}
\begin{itemize}
    \item \textbf{Gravimetría directa:} El analito se convierte en un compuesto poco soluble y se pesa.
    \item \textbf{Gravimetría indirecta:} Se mide la pérdida de masa de la muestra, por ejemplo, en un proceso de calcinación.
\end{itemize}

\section{Procedimiento experimental}
\begin{enumerate}
    \item \textbf{Secado y peso constante:} El material de laboratorio (ej. crisol) se seca en una estufa a 105°C hasta alcanzar un peso constante.
    \item \textbf{Precipitación:} Se añade un reactivo precipitante a la muestra para formar un precipitado.
    \item \textbf{Digestión:} El precipitado se calienta para mejorar la pureza del cristal.
    \item \textbf{Filtración y lavado:} Se separa el precipitado y se lava para eliminar impurezas.
    \item \textbf{Secado y calcinación:} Se seca o calcina el precipitado para obtener una forma de peso conocido.
    \item \textbf{Pesada:} Se pesa el precipitado y se calcula la cantidad de analito.
\end{enumerate}

\resumebox{ejercicio}{Cálculo gravimétrico}{\footnotesize
    Una muestra de 0.5 g de mineral de hierro se disuelve y se precipita como $Fe(OH)_3$, que se calcina a $Fe_2O_3$. El peso del $Fe_2O_3$ obtenido es de 0.2 g. Calcula el porcentaje de hierro en el mineral. (Masas atómicas: Fe=55.85, O=16.00)
}